% @see : https://coopmaths.fr/alea/?uuid=f1208&id=2N30-flash3&n=2&d=10&qcm=0&cd=1&tip=0&alea=ka4d&uuid=97008&id=2N30-flash2&n=1&d=10&qcm=0&cd=1&tip=0&alea=dpSN&uuid=16c8e&id=2N30-flash1&n=2&d=10&cd=1&tip=0&alea=nY7r&uuid=d309b&id=2N30-4&n=1&d=10&s=false&s2=10&s3=true&s4=true&cd=0&tip=0&alea=q00i&uuid=d0fdc&id=2N33-1&n=1&d=10&s=6&s2=false&cd=1&tip=0&alea=Mdan&uuid=08a2a&id=2N31-flash2&n=2&d=10&qcm=0&cd=1&tip=0&alea=w7p1&uuid=dd7d1&id=2N31-flash1&alea=PvQN&uuid=1894f&id=2N31-flash4&alea=C68e&uuid=07df0&id=2N31-flash3&alea=Ship&uuid=6575c&id=2N32-1&n=1&d=10&s=4&s2=true&cd=0&tip=0&alea=xvAD&uuid=4c878&id=2N33-flash1&alea=LIo0&v=eleve
\documentclass[a4paper,11pt,fleqn]{article}

%%% COULEURS %%%
\usepackage[table,svgnames]{xcolor}
\definecolor{nombres}{cmyk}{0,.8,.95,0}
\definecolor{gestion}{cmyk}{.75,1,.11,.12}
\definecolor{gestionbis}{cmyk}{.75,1,.11,.12}
\definecolor{grandeurs}{cmyk}{.02,.44,1,0}
\definecolor{geo}{cmyk}{.62,.1,0,0}
\definecolor{algo}{cmyk}{.69,.02,.36,0}
\definecolor{correction}{cmyk}{.63,.23,.93,.06}
\arrayrulecolor{couleur_theme} % Couleur des filets des tableaux
\usepackage{tasks}

\usepackage[most]{tcolorbox} % Supprimé à cause de conflits avec ProfCollege
% Découpe des cadres d'exercice. Le style est redéfini par la sortie LaTeX
% juste avant un exercice à garder d'un seul tenant, puis rétabli :
%   \tcbset{mathaleaexo/.style={unbreakable}}
% Passer par un style plutôt que par la clé « breakable » directement évite
% de changer la découpe des autres cadres du document (ProfCollege...).
\tcbset{mathaleaexo/.style={breakable}}
\tcbuselibrary{documentation}

%%%%%%%% Modification paramétrage pour footnotes
% Le paquet semble déjà appelé en amont, curieux qu'il n'y ait pas eu de clash
% avec l'appel ci-dessous mais seulement quand j'ai essayé de l'inclure avec l'option lualatex

% \usepackage{hyperref}
% Parametrages
\hypersetup{
    colorlinks=true,% On active la couleur pour les liens. Couleur par défaut rouge
    linkcolor=black,% On définit la couleur pour les liens internes
    % filecolor=magenta,% On définit la couleur pour les liens vers les fichiers locaux      
    urlcolor=black,% On définit la couleur pour les liens vers des sites web
    % pdftitle={Puissance Quatre},% On définit un titre pour le document pdf
    % pdfpagemode=FullScreen,% On fixe l'affichage par défaut à plein écran
    }
% Pour avoir les numérotations arabic y compris dans les environnements tcolorbox
\renewcommand{\thempfootnote}{\arabic{mpfootnote}}
%%%%%%%%

\usepackage{multicol} 
\usepackage{calc}
\usepackage{enumitem}
\usepackage{graphicx}
\usepackage{tabularx}
\usepackage{tablvar}

\setlength{\parindent}{0mm}
\renewcommand{\arraystretch}{1.5}
\renewcommand{\labelenumi}{\textbf{\theenumi{}.}}
\renewcommand{\labelenumii}{\textbf{\theenumii{}.}}
\setlength{\fboxsep}{3mm}

\setlength{\headheight}{14.5pt}

\setlength{\premulticols}{0pt}

\newcounter{ExoMA}

\newlength{\parindentMA}%
\setlength{\parindentMA}{\parindent}
\addtolength{\parindentMA}{30pt}

\usepackage{simplekv}

\setKVdefault[Theme]{Coopmath,Classiques=false,CANs=false}
\defKV[Theme]{Classique=\setKV[Theme]{Coopmath=false}\setKV[Theme]{Classiques}}
\defKV[Theme]{CAN=\setKV[Theme]{Coopmath=false}\setKV[Theme]{CANs}}

\usepackage{fancyhdr}
\fancyhead[C]{}
\fancyfoot{}

\def\bla{}

\NewDocumentCommand\Theme{o m m m m}{%
  \useKVdefault[Theme]%
  \setKV[Theme]{#1}%
  \ifboolKV[Theme]{CANs}{%
    \usepackage[left=1.5cm,right=1.5cm,top=2cm,bottom=2cm]{geometry}%
    \fancypagestyle{premierePage}{%
      \fancyhead[C]{\textsc{\ifx\bla#2\bla Course aux nombres\else#2\fi}}%
      \fancyfoot[R]{\scriptsize Coopmaths.fr -- CC-BY-SA}%
      \setlength{\headheight}{14.5pt}%
      \NewDocumentCommand\dureeCan{}{\ifx\bla#4\bla 9 minutes\else #4\fi}
      \NewDocumentCommand\titreSujetCan{}{\ifx\bla#3\bla\else\textbf{#3}\fi}
    }%
    \pagestyle{premierePage}
    \colorlet{couleur_theme}{black}%
    \colorlet{couleur_numerotation}{couleur_theme}%
    \let\EXO\EXOCan\let\endEXO\endEXOCan%
  }{%
    \renewcommand{\headrulewidth}{0pt}
    \renewcommand{\footrulewidth}{0pt}
    \pagestyle{fancy}
    \ifboolKV[Theme]{Classiques}{%
      \usepackage[left=0.65cm,right=0.65cm,top=2cm,bottom=1.5cm]{geometry}
      \let\EXO\EXOlibre\let\endEXO\endEXOlibre%
      \colorlet{couleur_theme}{black}%
      \colorlet{couleur_numerotation}{couleur_theme}%
      \fancyhead[C]{\bfseries #3}
      \fancyhead[L]{\bfseries #4}
      \fancyfoot[R]{#5}
    }{%
      \usepackage[left=1.5cm,right=1.5cm,top=4cm,bottom=2cm]{geometry}
      \theme{#2}{#3}{#4}{#5}%
      \let\EXO\EXOcoop\let\endEXO\endEXOcoop%
    }%
  }%
}%

\NewDocumentEnvironment{EXOCan}{m m}{%
}{}

\NewDocumentEnvironment{EXOcoop}{m m}{%
  \begin{tcolorbox}[%
    enhanced,
    mathaleaexo,
    colback=white,
    colframe=white,
    coltitle=black,
    title=\IfNoValueTF{#1}{}{\textmd{#1}},%
    overlay unbroken and first={%
      \IfNoValueTF{#2}{}{%
        \node[
        fill=white,
        anchor=east,
        xshift=10pt,
        text=black
        ]
        at (frame.north east)
        {\scriptsize #2};
      }
      \node[anchor=west,xshift=-20pt,yshift=-15pt] 
      at (frame.north west) {\stepcounter{ExoMA}\numb{\arabic{ExoMA}}};
    }
  ]
}{%
  \end{tcolorbox}
}

\NewDocumentEnvironment{EXOlibre}{m m}{%
  \begin{tcolorbox}[%
    enhanced,
    mathaleaexo,
    colback=white,
    colframe=white,
    coltitle=black,
    title=\stepcounter{ExoMA}\textbf{Exercice \arabic{ExoMA}},%
  ]
  \IfNoValueTF{#1}{}{#1\par}%
}{%
  \end{tcolorbox}%
}


%%% MISE EN PAGE %%%
\usepackage{setspace}
\usepackage{pgf}%
\usepackage{pgfplots}
\pgfplotsset{compat=1.18}
\usetikzlibrary{babel,arrows, arrows.meta,calc,fit,patterns,plotmarks,shapes.geometric,shapes.misc,shapes.symbols,shapes.arrows,shapes.callouts, shapes.multipart, shapes.gates.logic.US,shapes.gates.logic.IEC, er, automata,backgrounds,chains,topaths,trees,petri,mindmap,matrix, calendar, folding,fadings,through,positioning,scopes,decorations.fractals,decorations.shapes,decorations.text,decorations.pathmorphing,decorations.pathreplacing,decorations.footprints,decorations.markings,shadows}

\spaceskip=2\fontdimen2\font plus 3\fontdimen3\font minus3\fontdimen4\font\relax %Pour doubler l'espace entre les mots
\newcommand\numb[1]{ % Dessin autour du numéro d'exercice
  \begin{tikzpicture}[overlay,scale=.8]%,yshift=-.3cm
    \draw[fill=couleur_numerotation,couleur_numerotation](-.4,-0.1)rectangle($(-0.4,.8)+(2em,0)$);
    \draw (-.4,.8) node[white,anchor=north west,inner sep=0.5pt]{\bfseries EX}; 
    \draw[line width=.05cm,couleur_numerotation,fill=white] (0,0)--(.5,.5)--(1,0)--(.5,-.5)--cycle;
    \draw (0.5,0) node[anchor=center,couleur_numerotation]{\large\bfseries#1};
  \end{tikzpicture}
}

% echelle pour le dé
\def\globalscale {0.04}
% abscisse initiale pour les chevrons
\def\xini {3}

\newcommand\theme[4]{%
  \fancyhead[C]{%
    % Tracé du dé
    \begin{tikzpicture}[y=0.80pt, x=0.80pt, yscale=-\globalscale, xscale=\globalscale,remember picture, overlay,transform canvas={xshift=-0.5\linewidth,yshift=2.7cm},fill=couleur_theme]%,xshift=17cm,yshift=9.5cm
      %%%% Arc supérieur gauche%%%%
      \path[fill](523,1424)..controls(474,1413)and(404,1372)..(362,1333)..controls(322,1295)and(313,1272)..(331,1254)..controls(348,1236)and(369,1245)..(410,1283)..controls(458,1328)and(517,1356)..(575,1362)..controls(635,1368)and(646,1375)..(643,1404)..controls(641,1428)and(641,1428)..(596,1430)..controls(571,1431)and(538,1428)..(523,1424)--cycle;
      %%%% Dé face supérieur%%%%
      \path[fill](512,1272)..controls(490,1260)and(195,878)..(195,861)..controls(195,854)and(198,846)..(202,843)..controls(210,838)and(677,772)..(707,772)..controls(720,772)and(737,781)..(753,796)..controls(792,833)and(1057,1179)..(1057,1193)..controls(1057,1200)and(1053,1209)..(1048,1212)..controls(1038,1220)and(590,1283)..(551,1282)..controls(539,1282)and(521,1278)..(512,1272)--cycle;
      %%%% Dé faces gauche et droite%%%%
      \path[fill](1061,1167)..controls(1050,1158)and(978,1068)..(900,967)..controls(792,829)and(756,777)..(753,756)--(748,729)--(724,745)..controls(704,759)and(660,767)..(456,794)..controls(322,813)and(207,825)..(200,822)..controls(193,820)and(187,812)..(187,804)..controls(188,797)and(229,688)..(279,563)..controls(349,390)and(376,331)..(391,320)..controls(406,309)and(462,299)..(649,273)..controls(780,254)and(897,240)..(907,241)..controls(918,243)and(927,249)..(928,256)..controls(930,264)and(912,315)..(889,372)..controls(866,429)and(848,476)..(849,477)..controls(851,479)and(872,432)..(897,373)..controls(936,276)and(942,266)..(960,266)..controls(975,266)and(999,292)..(1089,408)..controls(1281,654)and(1290,666)..(1290,691)..controls(1290,720)and(1104,1175)..(1090,1180)..controls(1085,1182)and (1071,1176)..(1061,1167)--cycle;
      %%%% Arc inférieur bas%%%%
      \path[fill](1329,861)..controls(1316,848)and(1317,844)..(1339,788)..controls(1364,726)and(1367,654)..(1347,591)..controls(1330,539)and(1338,522)..(1375,526)..controls(1395,528)and(1400,533)..(1412,566)..controls(1432,624)and(1426,760)..(1401,821)..controls(1386,861)and(1380,866)..(1361,868)..controls(1348,870)and(1334,866)..(1329,861)--cycle;
      %%%% Arc inférieur gauche%%%%
      \path[fill](196,373)..controls(181,358)and(186,335)..(213,294)..controls(252,237)and(304,190)..(363,161)..controls(435,124)and(472,127)..(472,170)..controls(472,183)and(462,192)..(414,213)..controls(350,243)and(303,283)..(264,343)..controls(239,383)and(216,393)..(196,373)--cycle;
    \end{tikzpicture}
    \begin{tikzpicture}[line cap=round,line join=round,remember picture, overlay, shift={(current page.north west)},yshift=-8.5cm]
      \fill[fill=couleur_theme] (0,5) rectangle (21,6);
      \fill[fill=couleur_theme] (\xini,6)--(\xini+1.5,6)--(\xini+2.5,7)--(\xini+1.5,8)--(\xini,8)--(\xini+1,7)-- cycle;
      \fill[fill=couleur_theme] (\xini+2,6)--(\xini+2.5,6)--(\xini+3.5,7)--(\xini+2.5,8)--(\xini+2,8)--(\xini+3,7)-- cycle;  
      \fill[fill=couleur_theme] (\xini+3,6)--(\xini+3.5,6)--(\xini+4.5,7)--(\xini+3.5,8)--(\xini+3,8)--(\xini+4,7)-- cycle;   
      \node[color=white] at (10.5,5.5) {\LARGE \bfseries{ \MakeUppercase{#4}}};
    \end{tikzpicture}
    \begin{tikzpicture}[remember picture,overlay]
      \node[anchor=north east,inner sep=0pt] at ($(current page.north east)+(0,-.8cm)$) {};
      \node[anchor=west, fill=white, inner sep = 0pt] at ($(current page.north west)+(0.2,-2.3cm)$) {\footnotesize \bfseries{MathALÉA}};
    \end{tikzpicture}
    \begin{tikzpicture}[remember picture,overlay]
      \node[anchor=north east,inner sep=0pt] at ($(current page.north east)+(0,-.8cm)$) {};
      \node[anchor=east, fill=white] at ($(current page.north east)+(-2,-1.5cm)$) {\Huge \textcolor{couleur_theme}{\bfseries{\#}} \bfseries{#2} \textcolor{couleur_theme}{\bfseries \MakeUppercase{#3}}};
    \end{tikzpicture}
  }%
  \fancyfoot[R]{%
    \begin{tikzpicture}[remember picture,overlay]
      \node[anchor=south east] at ($(current page.south east)+(-2,0.25cm)$) {\scriptsize {\bfseries \href{https://coopmaths.fr/}{Coopmaths.fr} -- \href{http://creativecommons.fr/licences/}{CC-BY-SA}}};
    \end{tikzpicture}
    \begin{tikzpicture}[line cap=round,line join=round,remember picture, overlay,xscale=0.5,yscale=0.5, shift={(current page.south west)},xshift=35.7cm,yshift=-6cm]
      \fill[fill=couleur_theme] (\xini,6)--(\xini+1.5,6)--(\xini+2.5,7)--(\xini+1.5,8)--(\xini,8)--(\xini+1,7)-- cycle;
      \fill[fill=couleur_theme] (\xini+2,6)--(\xini+2.5,6)--(\xini+3.5,7)--(\xini+2.5,8)--(\xini+2,8)--(\xini+3,7)-- cycle;  
      \fill[fill=couleur_theme] (\xini+3,6)--(\xini+3.5,6)--(\xini+4.5,7)--(\xini+3.5,8)--(\xini+3,8)--(\xini+4,7)-- cycle;  
    \end{tikzpicture}
  }
  \fancyfoot[C]{}
  \colorlet{couleur_theme}{#1}
  \colorlet{couleur_numerotation}{couleur_theme}
  \def\iconeobjectif{icone-objectif-#1}
  \def\urliconeomethode{icone-methode-#1}
}%

\newcommand*\tikzfootMA{%
  \ifboolKV[Theme]{CANs}{
    \begin{tikzpicture}[remember picture,overlay,shift=(current page.north west)]
      \begin{scope}[x=\textwidth-\oddsidemargin,y=\textheight+5pt]
        \draw[correction,line width=4pt,dashed,dash pattern= on 10pt off 10pt,shift={(-\oddsidemargin,\topmargin)}] (0,0) rectangle (1,-1);
      \end{scope}
    \end{tikzpicture} 
  }{%
  \ifboolKV[Theme]{Classiques}{%
  \begin{tikzpicture}[remember picture,overlay,shift=(current page.south west)]
    \begin{scope}[x={(current page.south east)},y={(current page.north west)}]
      \draw[correction,line width=4pt,dashed,dash pattern= on 10pt off 10pt] ($(0,0)+(5mm,1cm)$) rectangle ($(1,1)+(-5mm,-1.75cm)$);
    \end{scope}
  \end{tikzpicture} 
  }{%
  \begin{tikzpicture}[remember picture,overlay,shift=(current page.south west)]
    \begin{scope}[x={(current page.south east)},y={(current page.north west)}]
      \draw[correction,line width=4pt,dashed,dash pattern= on 10pt off 10pt] ($(0,0)+(5mm,1.5cm)$) rectangle ($(1,1)+(-5mm,-3.75cm)$);
    \end{scope}
  \end{tikzpicture} 
  }%
  }
}

\newenvironment{Correction}{%
  \setcounter{ExoMA}{0}%
  \cfoot{\tikzfootMA}
}{}%

\newcommand\version[1]{
  \fancyhead[R]{
    \begin{tikzpicture}[remember picture,overlay]
    \node[anchor=north east,inner sep=0pt] at ($(current page.north east)+(-.5,-.5cm)$) {\large \textcolor{couleur_theme}{\bfseries V#1}};
    \end{tikzpicture}
  }
}

\usepackage[french]{babel}
\usepackage{icomma}
% loadPackagesFromContent

\newcounter{customfootnote}
\newcounter{tempreset}%
\newcounter{temp}

\makeatletter
\newcommand{\anote}[1]{%
  \refstepcounter{customfootnote}%
  \textsuperscript{\thecustomfootnote}%
  \expandafter\gdef\csname custom@footnote@\thecustomfootnote\endcsname{#1}%
}

\newcommand{\printcustomnotes}{%
  \par\noindent\rule{\linewidth}{0.4pt}\par
  \setcounter{temp}{1}%
  \whileboolexpr{
    test {\ifcsdef{custom@footnote@\thetemp}}
  }{%
    \noindent\textsuperscript{\thetemp}~\csname custom@footnote@\thetemp\endcsname\par
    \stepcounter{temp}%
  }%
}

\newcommand{\resetcustomnotes}{%
  \setcounter{customfootnote}{0}%
  \setcounter{tempreset}{1}%
  \whileboolexpr{
    test {\ifcsdef{custom@footnote@\thetempreset}}
  }{%
    \global\expandafter\let\csname custom@footnote@\thetempreset\endcsname\@undefined
    \stepcounter{tempreset}%
  }%
}
\makeatother
\newcommand{\e}{\mathrm{e}}
\usepackage{amsmath}
\usepackage{etoolbox}
\newbool{dys}
\setbool{dys}{false}
% ===== POLICES =====
% Le choix est le même en police standard et en police adaptée aux dys ;
% seuls la taille et l'interligne changent entre les deux.
\newcommand{\choiceFonts}[1]{
\ifstrequal{#1}{helvet}{%
  % Helvetica (police standard historique des fiches)
  \usepackage[scaled=1]{helvet}
}{}
\ifstrequal{#1}{Fira}{%
  % Fira Sans + Fira Math
  % Description : Fira Sans est une police moderne, et Fira Math est une version mathématique compatible qui maintient le style sans-serif.
  % Utilisation : Fonctionne bien avec XeLaTeX et LuaLaTeX.
  \usepackage{fontspec}
  \setmainfont{Fira Sans}
  \setsansfont{Fira Sans}
  \usepackage{unicode-math}
  \setmathfont{Fira Math}
}{}
\ifstrequal{#1}{lmodern}{%
  % Latin Modern Sans
  % Description : Une version sans-serif de la célèbre police Latin Modern, qui est compatible avec mathastext.
  % Utilisation : Fonctionne avec pdfLaTeX, XeLaTeX et LuaLaTeX.
  \usepackage{lmodern}
  \renewcommand{\familydefault}{\sfdefault}
  \usepackage{mathastext}
}{}
\ifstrequal{#1}{tgheros}{%
  % TeX Gyre Heros + mathastext
  % Description : TeX Gyre Heros est une alternative à Helvetica, et le package mathastext permet d'utiliser la même police pour les mathématiques.
  % Utilisation : Fonctionne avec pdfLaTeX, XeLaTeX, ou LuaLaTeX.
  \usepackage{tgheros}
  \renewcommand{\familydefault}{\sfdefault}
  \usepackage{mathastext}
}{}
}
\ifbool{dys}{
% POLICE DYS
\choiceFonts{Fira}

%\usepackage{unicode-math}
%\usepackage{fontspec}
%\setmainfont{TeX Gyre Schola}
%\setsansfont{TeX Gyre Schola}
%\setmathfont{TeX Gyre Schola Math}
%\setmainfont{OpenDyslexic}[Scale=1.0]
%\setsansfont{OpenDyslexic}[Scale=1.0]
%\setmathfont{OpenDyslexic}[Scale=1.0,range=up/{Latin,latin,num}]
\usepackage[fontsize=14]{scrextend}
\usepackage{setspace}
\setstretch{1.7}
% Une valeur d'environ 1.2 à 1.7 est souvent recommandée. Cela permet d'augmenter l'espace entre les lignes tout en offrant un léger espacement entre les mots.
\setlength{\spaceskip}{1.2em}
% Une valeur d'environ 1.2em à 1.5em est couramment conseillée. Cela crée un espace plus ample entre les mots, ce qui peut aider à réduire la fatigue visuelle et à améliorer la fluidité de la lecture.
}{
% POLICE STANDARD
%%% EE (24/04/2026) : Cette modif ci-dessous est nécessaire pour accepter ’ comme apostrophe.
% \usepackage[T1]{fontenc}     % Réservé à pdfLaTeX, à remplacer par fontspec en LuaLaTeX
\choiceFonts{helvet}
\usepackage[fontsize=11]{scrextend}
}

\Theme[Classique]{nombres}{Fractions 1}{Secondes}{Accompagnement personnalisé}
\begin{document}

\tcbset{mathaleaexo/.style={breakable}}
% @see : https://coopmaths.fr/alea?uuid=f1208&id=2N30-flash3&n=2&d=10&alea=ka4d&cd=1&cols=1
\begin{EXO}{}{}


\begin{multicols}{2}
\begin{enumerate}
	\item \begin{minipage}[t]{\linewidth} Rendre  la fraction $\dfrac{72}{81}$ irréductible. \end{minipage}
	\item \begin{minipage}[t]{\linewidth} Rendre  la fraction $\dfrac{50}{70}$ irréductible. \end{minipage}
\end{enumerate}
\end{multicols}
\tcbset{mathaleaexo/.style={breakable}}
% @see : https://coopmaths.fr/alea?uuid=97008&id=2N30-flash2&n=1&d=10&alea=dpSN&cd=1&cols=1
\end{EXO}

\begin{EXO}{}{}
\begin{enumerate}
    \item 
Écrire $\dfrac{70}{9}$ sous la forme de la somme d'un nombre entier et d'une fraction inférieure à 1.


    \item $\dfrac{70}{9} = \ldots + \ldots$
\end{enumerate}
\tcbset{mathaleaexo/.style={breakable}}
% @see : https://coopmaths.fr/alea?uuid=16c8e&id=2N30-flash1&n=2&d=10&alea=nY7r&cd=1&cols=1
\end{EXO}

\begin{EXO}{}{}
Compléter avec $>$ ou $<$.


\begin{tasks}(2)
	\task $\dfrac{17}{32}$ $\ldots$ $\dfrac{14}{5}$
	\task $\dfrac{13}{22}$ $\ldots$ $\dfrac{14}{5}$
\end{tasks}
\tcbset{mathaleaexo/.style={breakable}}
% @see : https://coopmaths.fr/alea?uuid=d309b&id=2N30-4&n=1&d=10&s=false&s2=10&s3=true&s4=true&alea=q00i&cd=0&cols=2
\end{EXO}

\begin{EXO}{
\resetcustomnotes }{}

Compléter avec deux nombres entiers \textbf{consécutifs} %\anote{\textbf{} Nombres entiers consécutifs : Ce sont deux nombres entiers qui se suivent comme 4 et 5.}.

 $$\ldots\ldots < \dfrac{5}{2} < \ldots\ldots$$
%\printcustomnotes
\tcbset{mathaleaexo/.style={breakable}}
% @see : https://coopmaths.fr/alea?uuid=d0fdc&id=2N33-1&n=1&d=10&s=6&s2=false&alea=Mdan&cd=1&cols=1
\end{EXO}

\begin{EXO}{}{}
Justifier votre réponse au problème suivant.


 Un jardin est aménagé selon les proportions suivantes :  $\dfrac{1}{5}$ par la culture des légumes, $\dfrac{19}{100}$ par  la culture des plantes aromatiques, $\dfrac{6}{25}$ par  une serre servant aux semis et le reste par la culture des fraisiers. 

 Quelle est la culture qui occupe le plus de surface ?


\tcbset{mathaleaexo/.style={breakable}}
% @see : https://coopmaths.fr/alea?uuid=08a2a&id=2N31-flash2&n=2&d=10&alea=w7p1&cd=1&cols=1
\end{EXO}

\begin{EXO}{}{}


\begin{tasks}(3)
	\task \begin{minipage}[t]{\linewidth} Calculer $\dfrac{3}{4} + \dfrac{2}{3}$. \end{minipage}
	\task \begin{minipage}[t]{\linewidth} Calculer $\dfrac{8}{9} - \dfrac{2}{7}$. \end{minipage}
    \task Calculer $3 - \dfrac{5}{8}$.
\end{tasks}
\tcbset{mathaleaexo/.style={breakable}}
% @see : https://coopmaths.fr/alea?uuid=dd7d1&id=2N31-flash1&n=1&d=10&alea=PvQN&cd=1&cols=1
% @see : https://coopmaths.fr/alea?uuid=1894f&id=2N31-flash4&n=1&d=10&alea=C68e&cd=1&cols=1
\end{EXO}

\begin{EXO}{}{}
\begin{enumerate}

\item Calculer  $\dfrac{13}{20}\times \dfrac{36}{13}$ sous la forme d'une fraction irréductible. 
% @see : https://coopmaths.fr/alea?uuid=07df0&id=2N31-flash3&n=1&d=10&alea=Ship&cd=1&cols=1



\item Calculer et écrire sous la forme d'une fraction simplifiée : $\dfrac{3}{4}\times \dfrac{7}{9}$.


% @see : https://coopmaths.fr/alea?uuid=6575c&id=2N32-1&n=1&d=10&s=4&s2=true&alea=xvAD&cd=0&cols=4

\item Effectuer le calcul suivant en donnant le résultat sous forme d'une fraction.

 $$A = \dfrac{9}{7} \times \dfrac{4}{7} - \dfrac{2}{7}$$
\end{enumerate}
\tcbset{mathaleaexo/.style={breakable}}
% @see : https://coopmaths.fr/alea?uuid=4c878&id=2N33-flash1&n=1&d=10&alea=LIo0&cd=1&cols=1
\end{EXO}

\begin{EXO}{}{}



 Ce matin, Alice a ouvert une bouteille d’eau.

     Elle a bu $\dfrac{1}{3}$ de la bouteille. Puis à midi,  elle a bu $\dfrac{3}{4}$ du reste.

 
Quelle fraction de la bouteille a-t-elle bu à midi ? 
\end{EXO}


\clearpage

\begin{Correction}
\begin{EXO}{}{}

\begin{enumerate}
	\item $\dfrac{72}{81} =\dfrac{8{\color[HTML]{2563a5}\boldsymbol{\times9}} }{9{\color[HTML]{2563a5}\boldsymbol{\times9}}}={\color[HTML]{F15929}\boldsymbol{\dfrac{8}{9}}}$
	\item $\dfrac{50}{70} =\dfrac{5{\color[HTML]{2563a5}\boldsymbol{\times10}} }{7{\color[HTML]{2563a5}\boldsymbol{\times10}}}={\color[HTML]{F15929}\boldsymbol{\dfrac{5}{7}}}$
\end{enumerate}
\end{EXO}

\begin{EXO}{}{}

 Le plus grand multiple de $9$ inférieur à $70$ est $63$. 

    Ainsi, $\dfrac{70}{9}=\dfrac{63}{9}+\dfrac{7}{9}={\color[HTML]{F15929}\boldsymbol{7+\dfrac{7}{9}}}$.
\end{EXO}

\begin{EXO}{}{}

\begin{enumerate}
	\item $\dfrac{17}{32} <1$ et $\dfrac{14}{5}>1$.

      On en déduit : $\dfrac{17}{32} {\color[HTML]{F15929}\boldsymbol{<}} \dfrac{14}{5}$.
	\item $\dfrac{13}{22} <1$ et $\dfrac{14}{5}>1$.

      On en déduit : $\dfrac{13}{22} {\color[HTML]{F15929}\boldsymbol{<}} \dfrac{14}{5}$.
\end{enumerate}
\end{EXO}

\begin{EXO}{}{}

  $\quad \dfrac{5}{2}=2+\dfrac{1}{2}\quad$ et $\quad2<2+\dfrac{1}{2}<3$ 

donc $\quad {\color[HTML]{F15929}\boldsymbol{2}} < \dfrac{5}{2} < {\color[HTML]{F15929}\boldsymbol{3}}$.
\end{EXO}

\begin{EXO}{}{}

 Il s'agit d'un problème additif. Il va être nécessaire de réduire les fractions au même dénominateur pour les additionner, les soustraire ou les comparer.

Réduisons les fractions de l'énoncé au même dénominateur :  $\dfrac{1}{5}$ $= \dfrac{20}{100}$ ; $\dfrac{19}{100}$  et $\dfrac{6}{25}$ $= \dfrac{24}{100}$.

Calculons d'abord la fraction du jardin occupée par une serre servant aux semis : 

$1-\dfrac{1}{5}-\dfrac{19}{100}-\dfrac{6}{25} = \dfrac{100}{100}-\dfrac{20}{100}-\dfrac{19}{100}-\dfrac{24}{100} = \dfrac{100-20-19-24}{100} = \dfrac{37}{100}$

Le jardin est donc occupé de la façon suivante : $\dfrac{1}{5}$ par la culture des légumes, $\dfrac{19}{100}$ par la culture des plantes aromatiques, $\dfrac{6}{25}$ par une serre servant aux semis et $\dfrac{37}{100}$ par la culture des fraisiers.

 Avec les mêmes dénominateurs pour pouvoir comparer, le jardin est donc occupé de la façon suivante : $\dfrac{20}{100}$ par la culture des légumes, $\dfrac{19}{100}$ par la culture des plantes aromatiques, $\dfrac{24}{100}$ par une serre servant aux semis et $\dfrac{37}{100}$ par la culture des fraisiers.

Nous pouvons alors ranger ces fractions dans l'ordre croissant : $\dfrac{19}{100}$, $\dfrac{20}{100}$, $\dfrac{24}{100}$, $\dfrac{37}{100}$.

Enfin, nous pouvons ranger les fractions de l'énoncé et la fraction calculée dans l'ordre croissant : ${\color[HTML]{F15929}\boldsymbol{\dfrac{19}{100}}} < {\color[HTML]{F15929}\boldsymbol{\dfrac{1}{5}}} < {\color[HTML]{F15929}\boldsymbol{\dfrac{6}{25}}} < {\color[HTML]{F15929}\boldsymbol{\dfrac{37}{100}}}$.

 
                      C'est donc par {\bfseries \color[HTML]{F15929}la culture des fraisiers} que le jardin est le plus occupé.
\end{EXO}

\begin{EXO}{}{}

\begin{enumerate}[itemsep=1em]
	\item On cherche un dénominateur commun (ici $12$), puis on additionne les numérateurs :

$\begin{aligned}
\dfrac{3}{4} + \dfrac{2}{3} &= \dfrac{3\times 3}{4\times 3} + \dfrac{2\times 4}{3\times 4} \\
&= \dfrac{9}{12} + \dfrac{8}{12} \\[0.7em]
&= {\color[HTML]{F15929}\boldsymbol{\dfrac{17}{12}}}
\end{aligned}$
	\item On réduit au même dénominateur ($63$), puis on soustrait les numérateurs :

$\begin{aligned}
\dfrac{8}{9} - \dfrac{2}{7} &= \dfrac{8\times 7}{9\times 7} -\dfrac{2\times 9}{7\times 9} \\
&= \dfrac{56}{63} - \dfrac{18}{63} \\[0.7em]
&= {\color[HTML]{F15929}\boldsymbol{\dfrac{38}{63}}}
\end{aligned}$
\end{enumerate}
\medskip


 On écrit l'entier sous la forme d'une fraction de dénominateur $8$, soit $3=  \dfrac{24}{8}$

$\begin{aligned}
3 - \dfrac{5}{8} &= \dfrac{24}{8} - \dfrac{5}{8} \\
&= {\color[HTML]{F15929}\boldsymbol{\dfrac{19}{8}}}
\end{aligned}$
\end{EXO}

\begin{EXO}{}{}

 Il y a une simplification évidente par $13$, puis on simplifie par $4$ :

    $\begin{aligned}
    \dfrac{13}{20}\times \dfrac{36}{13}&=\dfrac{36}{20}\\
    &=\dfrac{4 \times 9}{4 \times 5}\\
    &={\color[HTML]{F15929}\boldsymbol{\dfrac{9}{5}}}
    \end{aligned}$
\medskip


 $\dfrac{3}{4}\times \dfrac{7}{9}=\dfrac{3\times 7}{4 \times 9}=\dfrac{3\times 7}{2\times2\times 3\times3}=\dfrac{7{\color[HTML]{2563a5}\boldsymbol{\times3}} }{12{\color[HTML]{2563a5}\boldsymbol{\times3}}}=\dfrac{7}{12}$
\medskip

\begin{multicols}{4}
 \begin{minipage}[t]{\linewidth} $A = \dfrac{9}{7} \times \dfrac{4}{7} - \dfrac{2}{7}$

$A = \dfrac{36}{49} - \dfrac{2}{7}$

$A = \dfrac{36}{49} - \dfrac{14}{49}$

$A =$ ${\color[HTML]{F15929}\boldsymbol{ \dfrac{22}{49}}}$

 \end{minipage}
\end{multicols}
\end{EXO}

\begin{EXO}{}{}

 Le matin, Alice a bu $\dfrac{1}{3}$ de la bouteille. Il reste alors $\dfrac{2}{3}$ de la bouteille.

    À midi, elle a bu $\dfrac{3}{4}$ du reste.

     Comme $\dfrac{3}{4}\times \dfrac{2}{3}=\dfrac{6}{12}$, elle a bu ${\color[HTML]{F15929}\boldsymbol{\dfrac{6}{12}}}$ 
      ou $\dfrac{1}{2}$ de la bouteille à midi.
\end{EXO}

\clearpage
\end{Correction}
\end{document}

% Local Variables:
% TeX-engine: luatex
% End: